\chapter{Probability and Statistics}
\label{ch:probability}
% ============================================================================

Probability theory is essential for stochastic control, Kalman filtering, and machine learning. This chapter provides the mathematical foundations needed to understand uncertainty quantification, state estimation, and statistical inference in control systems.

\section{Probability Fundamentals}

\begin{definitionbox}[title=Definition: Probability Space]
A \textbf{probability space} is a triple $(\Omega, \mathcal{F}, P)$ where:
\begin{itemize}
    \item $\Omega$ is the \textbf{sample space} (set of all possible outcomes)
    \item $\mathcal{F}$ is the \textbf{$\sigma$-algebra} (collection of events)
    \item $P: \mathcal{F} \to [0,1]$ is the \textbf{probability measure}
\end{itemize}
\end{definitionbox}

\subsubsection{Axioms of Probability}
\begin{enumerate}
    \item $P(A) \geq 0$ for all events $A$
    \item $P(\Omega) = 1$ where $\Omega$ is the sample space
    \item For mutually exclusive events: $P(A \cup B) = P(A) + P(B)$
\end{enumerate}

\subsubsection{Conditional Probability}

\begin{definitionbox}[title=Definition: Conditional Probability]
The \textbf{conditional probability} of event $A$ given event $B$ (with $P(B) > 0$) is:
\begin{equation}
P(A|B) = \frac{P(A \cap B)}{P(B)}
\end{equation}
\end{definitionbox}

\subsubsection{Bayes' Theorem}

\begin{keyconceptbox}[title=Bayes' Theorem]
For events $A$ and $B$ with $P(B) > 0$:
\begin{equation}
P(A|B) = \frac{P(B|A)P(A)}{P(B)}
\label{eq:bayes_theorem}
\end{equation}
In the context of estimation:
\begin{itemize}
    \item $P(A)$ is the \textbf{prior} probability
    \item $P(B|A)$ is the \textbf{likelihood}
    \item $P(A|B)$ is the \textbf{posterior} probability
    \item $P(B)$ is the \textbf{evidence} (normalizing constant)
\end{itemize}
\end{keyconceptbox}

The denominator can be expanded using the law of total probability:
\begin{equation}
P(B) = \sum_i P(B|A_i)P(A_i)
\end{equation}
where $\{A_i\}$ forms a partition of the sample space.

\begin{figure}[htbp]
\centering
\begin{tikzpicture}

% --- Top Row: Three Distribution Plots ---

% Prior Distribution
\begin{axis}[
    name=prior,
    width=5.5cm, height=5cm,
    xlabel={Parameter $\theta$},
    ylabel={Probability Density},
    xmin=-2, xmax=6,
    ymin=0, ymax=0.45,
    grid=major,
    grid style={gray!30},
    axis on top,
    title={\bfseries\color{UofSCAtlantic} Prior $p(\theta)$},
    title style={at={(0.5,1.0)}, anchor=south}
]

% Prior: Gaussian centered at 1 with wide variance
\addplot[fill=UofSCAtlantic!30, draw=UofSCAtlantic, very thick, domain=-2:6, samples=150]
    {1/(1.5*sqrt(2*pi)) * exp(-(x-1)^2/(2*1.5^2))} \closedcycle;

% Mean marker
\draw[UofSCAtlantic, thick, dashed] (1, 0) -- (1, 0.28);
\node[UofSCAtlantic, font=\footnotesize, above] at (1, 0.28) {$\mu_0=1$};

% Variance indicator
\draw[UofSCAtlantic, thick, |<->|] (1, 0.12) -- (2.5, 0.12);
\node[UofSCAtlantic, font=\footnotesize, above] at (1.75, 0.12) {$\sigma_0$};

\end{axis}

% Likelihood Function
\begin{axis}[
    name=likelihood,
    at={(prior.south east)},
    anchor=south west,
    xshift=0.8cm,
    width=5.5cm, height=5cm,
    xlabel={Parameter $\theta$},
    ylabel={},
    xmin=-2, xmax=6,
    ymin=0, ymax=0.45,
    grid=major,
    grid style={gray!30},
    axis on top,
    title={\bfseries\color{UofSCHorseshoe} Likelihood $p(x|\theta)$},
    title style={at={(0.5,1.0)}, anchor=south}
]

% Likelihood: Gaussian centered at 3.5 (where data suggests theta should be)
\addplot[fill=UofSCHorseshoe!30, draw=UofSCHorseshoe, very thick, domain=-2:6, samples=150]
    {1/(1.0*sqrt(2*pi)) * exp(-(x-3.5)^2/(2*1.0^2))} \closedcycle;

% Peak marker
\draw[UofSCHorseshoe, thick, dashed] (3.5, 0) -- (3.5, 0.42);
\node[UofSCHorseshoe, font=\footnotesize, above] at (3.5, 0.42) {$x_{obs}=3.5$};

% Annotation
\node[UofSCHorseshoe, font=\footnotesize, align=center] at (0.5, 0.32) {Data\\Evidence};

\end{axis}

% Posterior Distribution
\begin{axis}[
    name=posterior,
    at={(likelihood.south east)},
    anchor=south west,
    xshift=0.8cm,
    width=5.5cm, height=5cm,
    xlabel={Parameter $\theta$},
    ylabel={},
    xmin=-2, xmax=6,
    ymin=0, ymax=0.45,
    grid=major,
    grid style={gray!30},
    axis on top,
    title={\bfseries\color{UofSCGarnet} Posterior $p(\theta|x)$},
    title style={at={(0.5,1.0)}, anchor=south}
]

% Posterior: Narrower Gaussian between prior and likelihood means
% Using Bayesian update: posterior mean = (sigma_1^2 * mu_0 + sigma_0^2 * x) / (sigma_0^2 + sigma_1^2)
% posterior_mean = (1*1 + 2.25*3.5) / (2.25 + 1) = 8.875/3.25 = 2.73
% posterior_var = 1/(1/2.25 + 1/1) = 1/1.44 = 0.69, sigma = 0.83
\addplot[fill=UofSCGarnet!30, draw=UofSCGarnet, very thick, domain=-2:6, samples=150]
    {1/(0.83*sqrt(2*pi)) * exp(-(x-2.73)^2/(2*0.83^2))} \closedcycle;

% Posterior mean marker
\draw[UofSCGarnet, thick, dashed] (2.73, 0) -- (2.73, 0.49);
\node[UofSCGarnet, font=\footnotesize, above] at (2.73, 0.49) {$\mu_{post}$};

% Show it's between prior and likelihood
\draw[UofSCCongaree, thin, dotted] (1, 0.35) -- (1, 0.38);
\draw[UofSCCongaree, thin, dotted] (3.5, 0.35) -- (3.5, 0.38);
\draw[UofSCCongaree, thin, <->] (1, 0.36) -- (3.5, 0.36);
\node[UofSCCongaree, font=\tiny, above] at (2.25, 0.36) {Compromise};

\end{axis}

% --- Bottom Row: Process Flow Diagram ---
\begin{scope}[yshift=-6.5cm]

% Flow boxes
\node[draw=UofSCAtlantic, fill=UofSCAtlantic!15, minimum width=3cm, minimum height=1.2cm,
      align=center, font=\small, sharp corners] (priorbox) at (1.5, 0)
      {\textbf{Prior Belief}\\$p(\theta)$};

\node[draw=UofSCHorseshoe, fill=UofSCHorseshoe!15, minimum width=3cm, minimum height=1.2cm,
      align=center, font=\small, sharp corners] (databox) at (6.5, 0)
      {\textbf{New Data}\\$x_{obs}$};

\node[draw=UofSCGarnet, fill=UofSCGarnet!15, minimum width=3.5cm, minimum height=1.2cm,
      align=center, font=\small, sharp corners] (posteriorbox) at (12, 0)
      {\textbf{Updated Belief}\\$p(\theta|x)$};

% Multiplication node
\node[draw=UofSCBlack, fill=white, circle, minimum size=0.8cm, font=\large] (mult) at (4, 0) {$\times$};

% Proportionality node
\node[draw=UofSCBlack, fill=white, minimum width=1.5cm, minimum height=0.8cm,
      font=\small, sharp corners] (prop) at (9, 0) {$\propto$};

% Arrows
\draw[->, thick, UofSCBlack] (priorbox.east) -- (mult.west);
\draw[->, thick, UofSCBlack] (databox.west) -- (mult.east);
\draw[->, thick, UofSCBlack] (mult.east) ++(0.3,0) -- (prop.west);
\draw[->, thick, UofSCBlack] (prop.east) -- (posteriorbox.west);

% Likelihood label above data arrow
\node[UofSCHorseshoe, font=\footnotesize] at (5.25, 0.5) {Likelihood $p(x|\theta)$};

% Bayes' Theorem equation
\node[font=\normalsize, align=center] at (6.75, -1.5)
    {$\displaystyle p(\theta|x) = \frac{p(x|\theta) \cdot p(\theta)}{p(x)}$};
\node[font=\small, UofSC70Black] at (6.75, -2.2)
    {Posterior $\propto$ Likelihood $\times$ Prior};

% Iteration arrow (feedback)
\draw[->, thick, UofSCRose, dashed] (posteriorbox.south) -- ++(0, -0.7) -| (priorbox.south)
    node[pos=0.25, below, font=\footnotesize, UofSCRose] {Next iteration: posterior becomes new prior};

\end{scope}

\end{tikzpicture}
\caption{Bayesian updating process showing how prior beliefs are combined with data to form posterior estimates. \textbf{Top row}: The prior distribution (Atlantic blue, $\mu_0=1$, broad uncertainty) represents initial knowledge before observing data. The likelihood function (Horseshoe green) quantifies how probable the observed data $x_{obs}=3.5$ is for each possible parameter value. The posterior distribution (Garnet) is the normalized product, representing updated knowledge---narrower than the prior (reduced uncertainty) with mean shifted toward the data. \textbf{Bottom}: Schematic of Bayes' theorem. The dashed Rose arrow shows sequential updating: each posterior becomes the prior for the next observation. This framework underlies Kalman filtering, recursive estimation, and probabilistic inference in control systems.}
\label{fig:bayes_update}
\end{figure}


\begin{examplebox}[title=Example: Computing Conditional Probability with Bayes' Theorem]
A sensor for detecting system faults has the following characteristics:
\begin{itemize}
    \item The sensor correctly detects a fault (true positive) 95\% of the time: $P(\text{alarm}|\text{fault}) = 0.95$
    \item The sensor has a false alarm rate of 3\%: $P(\text{alarm}|\text{no fault}) = 0.03$
    \item From historical data, the probability of a fault occurring is 1\%: $P(\text{fault}) = 0.01$
\end{itemize}

\textbf{Question:} If the alarm triggers, what is the probability that an actual fault occurred?

\textbf{Solution:}

We seek $P(\text{fault}|\text{alarm})$. Using Bayes' theorem:
\begin{equation*}
P(\text{fault}|\text{alarm}) = \frac{P(\text{alarm}|\text{fault})P(\text{fault})}{P(\text{alarm})}
\end{equation*}

First, compute $P(\text{alarm})$ using total probability:
\begin{align*}
P(\text{alarm}) &= P(\text{alarm}|\text{fault})P(\text{fault}) + P(\text{alarm}|\text{no fault})P(\text{no fault}) \\
&= (0.95)(0.01) + (0.03)(0.99) \\
&= 0.0095 + 0.0297 = 0.0392
\end{align*}

Now apply Bayes' theorem:
\begin{equation*}
P(\text{fault}|\text{alarm}) = \frac{(0.95)(0.01)}{0.0392} = \frac{0.0095}{0.0392} \approx 0.242
\end{equation*}

\textbf{Interpretation:} Despite the sensor having a 95\% detection rate, when an alarm triggers, there is only about a 24\% chance of an actual fault. This counterintuitive result (the \textit{base rate fallacy}) occurs because faults are rare events. The false alarms from the large population of non-fault cases dominate.

\textbf{Engineering implication:} This analysis is critical for designing alert systems in control applications. A second confirmation measurement or a more stringent threshold may be needed to reduce false positives.
\end{examplebox}

\section{Random Variables}

\subsubsection{Discrete Random Variables}
\begin{itemize}
    \item Probability mass function: $p(x) = P(X = x)$
    \item Expected value: $E[X] = \sum_x x \cdot p(x)$
    \item Variance: $\text{Var}(X) = E[(X - E[X])^2] = E[X^2] - (E[X])^2$
\end{itemize}

\subsubsection{Continuous Random Variables}
\begin{itemize}
    \item Probability density function: $f(x)$ where $P(a \leq X \leq b) = \int_a^b f(x)\,\dd x$
    \item Expected value: $E[X] = \int_{-\infty}^{\infty} x f(x)\,\dd x$
    \item Cumulative distribution: $F(x) = P(X \leq x) = \int_{-\infty}^x f(t)\,\dd t$
\end{itemize}

\begin{keyconceptbox}[title=Properties of Expectation and Variance]
For random variables $X$ and $Y$, and constants $a, b$:
\begin{align}
E[aX + b] &= aE[X] + b \\
\text{Var}(aX + b) &= a^2\text{Var}(X) \\
E[X + Y] &= E[X] + E[Y] \quad \text{(always true)}
\end{align}
If $X$ and $Y$ are independent:
\begin{equation}
\text{Var}(X + Y) = \text{Var}(X) + \text{Var}(Y)
\end{equation}
\end{keyconceptbox}

\section{Common Distributions}

\begin{figure}[htbp]
\centering
\begin{tikzpicture}
\begin{axis}[
    width=13cm, height=8cm,
    xlabel={$x$},
    ylabel={Probability Density $f(x)$},
    xmin=-4, xmax=6,
    ymin=0, ymax=0.55,
    legend pos=north east,
    grid=major,
    grid style={gray!30},
    axis on top,
    major tick length=0.08cm
]

% Gaussian distribution: mu=0, sigma=1
\addplot[UofSCGarnet, very thick, domain=-4:6, samples=200]
    {1/(sqrt(2*pi)) * exp(-x^2/2)};
\addlegendentry{Gaussian: $\mathcal{N}(0, 1)$}

% Uniform distribution: a=1, b=4
\addplot[UofSCAtlantic, very thick, domain=-4:6, samples=200]
    {(x >= 1 && x <= 4) ? 1/3 : 0};
\addlegendentry{Uniform: $U(1, 4)$}

% Exponential distribution: lambda=0.7
\addplot[UofSCHorseshoe, very thick, domain=0:6, samples=200]
    {0.7 * exp(-0.7*x)};
\addlegendentry{Exponential: $\lambda = 0.7$}

% Gaussian with different parameters: mu=2, sigma=0.8
\addplot[UofSCCongaree, very thick, dashed, domain=-4:6, samples=200]
    {1/(0.8*sqrt(2*pi)) * exp(-(x-2)^2/(2*0.8^2))};
\addlegendentry{Gaussian: $\mathcal{N}(2, 0.64)$}

% Mark the means
\draw[UofSCGarnet, thin, dashed] (0, 0) -- (0, 0.4);
\node[UofSCGarnet, font=\footnotesize, below] at (0, 0.02) {$\mu=0$};

\draw[UofSCAtlantic, thin, dashed] (2.5, 0) -- (2.5, 0.35);
\node[UofSCAtlantic, font=\footnotesize, below] at (2.5, 0.02) {$\bar{x}=2.5$};

% Annotation for exponential decay
\node[UofSCHorseshoe, font=\footnotesize, right] at (3.5, 0.2) {$f(x) = \lambda\ee^{-\lambda x}$};
\draw[UofSCHorseshoe, thin, ->] (3.3, 0.18) -- (2.0, 0.08);

% Annotation for Gaussian shape
\node[UofSCGarnet, font=\footnotesize, left] at (-2, 0.35) {Bell curve};
\draw[UofSCGarnet, thin, ->] (-1.8, 0.33) -- (-0.8, 0.28);

\end{axis}
\end{tikzpicture}
\caption{Comparison of common probability distributions used in control systems and stochastic modeling. The Gaussian (normal) distribution (Garnet, solid) centered at $\mu=0$ with $\sigma=1$ is the most prevalent in estimation and filtering. The Uniform distribution (Atlantic blue) models bounded uncertainty with constant probability over $[1, 4]$. The Exponential distribution (Horseshoe green) characterizes waiting times and failure rates. A second Gaussian (Congaree teal, dashed) shows how shifting the mean and adjusting variance changes the distribution shape. Note how each distribution integrates to unity over its support.}
\label{fig:probability_distributions}
\end{figure}


\subsubsection{Gaussian (Normal) Distribution}
\begin{equation}
f(x) = \frac{1}{\sigma\sqrt{2\pi}} \exp\left(-\frac{(x-\mu)^2}{2\sigma^2}\right)
\end{equation}
Notation: $X \sim \mathcal{N}(\mu, \sigma^2)$

\begin{figure}[htbp]
\centering
\begin{tikzpicture}
\begin{axis}[
    width=14cm, height=9cm,
    xlabel={$x$},
    ylabel={Probability Density $f(x)$},
    xmin=-4.5, xmax=4.5,
    ymin=0, ymax=0.5,
    xtick={-4,-3,-2,-1,0,1,2,3,4},
    xticklabels={$\mu-4\sigma$,$\mu-3\sigma$,$\mu-2\sigma$,$\mu-\sigma$,$\mu$,$\mu+\sigma$,$\mu+2\sigma$,$\mu+3\sigma$,$\mu+4\sigma$},
    x tick label style={font=\small},
    grid=major,
    grid style={gray!30},
    axis on top,
    major tick length=0.08cm,
    clip=false
]

% Fill regions for 68-95-99.7 rule (from widest to narrowest)
% 99.7% region (3 sigma)
\addplot[fill=UofSCHorseshoe!25, draw=none, domain=-3:3, samples=100]
    {1/(sqrt(2*pi)) * exp(-x^2/2)} \closedcycle;

% 95% region (2 sigma)
\addplot[fill=UofSCAtlantic!35, draw=none, domain=-2:2, samples=100]
    {1/(sqrt(2*pi)) * exp(-x^2/2)} \closedcycle;

% 68% region (1 sigma)
\addplot[fill=UofSCGarnet!40, draw=none, domain=-1:1, samples=100]
    {1/(sqrt(2*pi)) * exp(-x^2/2)} \closedcycle;

% Main Gaussian curve
\addplot[UofSCBlack, very thick, domain=-4.5:4.5, samples=200]
    {1/(sqrt(2*pi)) * exp(-x^2/2)};

% Vertical lines at sigma boundaries
\draw[UofSCGarnet, thick, dashed] (-1, 0) -- (-1, 0.242);
\draw[UofSCGarnet, thick, dashed] (1, 0) -- (1, 0.242);

\draw[UofSCAtlantic, thick, dashed] (-2, 0) -- (-2, 0.054);
\draw[UofSCAtlantic, thick, dashed] (2, 0) -- (2, 0.054);

\draw[UofSCHorseshoe, thick, dashed] (-3, 0) -- (-3, 0.0044);
\draw[UofSCHorseshoe, thick, dashed] (3, 0) -- (3, 0.0044);

% Mean line
\draw[UofSCBlack, thick] (0, 0) -- (0, 0.42);
\node[UofSCBlack, font=\bfseries, above] at (0, 0.42) {Mean ($\mu$)};

% 68% label
\draw[UofSCGarnet, thick, <->] (-1, 0.32) -- (1, 0.32);
\node[UofSCGarnet, font=\bfseries, above] at (0, 0.32) {68.27\%};

% 95% label
\draw[UofSCAtlantic, thick, <->] (-2, 0.42) -- (2, 0.42);
\node[UofSCAtlantic, font=\bfseries, above] at (0, 0.42) {95.45\%};

% 99.7% label
\draw[UofSCHorseshoe, thick, <->] (-3, 0.48) -- (3, 0.48);
\node[UofSCHorseshoe, font=\bfseries, above] at (0, 0.48) {99.73\%};

% Tail probability annotations
\node[UofSCCongaree, font=\footnotesize, align=center] at (-3.8, 0.12) {Tail\\0.135\%};
\node[UofSCCongaree, font=\footnotesize, align=center] at (3.8, 0.12) {Tail\\0.135\%};

% PDF formula
\node[UofSCBlack, font=\small, align=left] at (-3.3, 0.38) {$f(x) = \frac{1}{\sigma\sqrt{2\pi}}\exp\left(-\frac{(x-\mu)^2}{2\sigma^2}\right)$};

% Standard deviation arrows
\draw[UofSCRose, thick, |<->|] (0, -0.03) -- (1, -0.03);
\node[UofSCRose, font=\footnotesize, below] at (0.5, -0.03) {$\sigma$};

\end{axis}
\end{tikzpicture}
\caption{The Gaussian (normal) distribution illustrating the 68-95-99.7 empirical rule. The probability density function is shown with shaded regions indicating the probability mass within one (Garnet, 68.27\%), two (Atlantic blue, 95.45\%), and three (Horseshoe green, 99.73\%) standard deviations of the mean. This rule is fundamental in statistical quality control, confidence interval estimation, and Kalman filter design. In control systems, sensor noise is often modeled as Gaussian with known $\mu$ and $\sigma$, making these probability regions essential for uncertainty quantification.}
\label{fig:normal_curve}
\end{figure}


\subsubsection{Multivariate Gaussian}
\begin{equation}
f(\vecx) = \frac{1}{(2\pi)^{n/2}|\boldsymbol{\Sigma}|^{1/2}} \exp\left(-\frac{1}{2}(\vecx-\boldsymbol{\mu})^\top\boldsymbol{\Sigma}^{-1}(\vecx-\boldsymbol{\mu})\right)
\end{equation}
where $\boldsymbol{\mu}$ is the mean vector and $\boldsymbol{\Sigma}$ is the covariance matrix.

\subsubsection{Other Important Distributions}
\begin{itemize}
    \item \textbf{Uniform:} $f(x) = \frac{1}{b-a}$ for $x \in [a,b]$
    \item \textbf{Exponential:} $f(x) = \lambda\ee^{-\lambda x}$ for $x \geq 0$
    \item \textbf{Bernoulli:} $P(X=1) = p$, $P(X=0) = 1-p$
    \item \textbf{Poisson:} $P(X=k) = \frac{\lambda^k \ee^{-\lambda}}{k!}$
\end{itemize}

\begin{examplebox}[title=Example: Mean and Variance of a Mixture of Distributions]
A manufacturing process has two operational modes. In Mode A (70\% of the time), output follows $\mathcal{N}(10, 1)$. In Mode B (30\% of the time), output follows $\mathcal{N}(15, 4)$. Find the overall mean and variance.

\textbf{Solution:}

Let $X$ be the output and $M$ be the mode indicator. The mixture distribution has PDF:
\begin{equation*}
f_X(x) = 0.7 \cdot f_A(x) + 0.3 \cdot f_B(x)
\end{equation*}

\textbf{Mean:} Using the law of total expectation:
\begin{align*}
E[X] &= E[E[X|M]] = P(M=A)E[X|M=A] + P(M=B)E[X|M=B] \\
&= 0.7 \times 10 + 0.3 \times 15 = 7 + 4.5 = 11.5
\end{align*}

\textbf{Variance:} Using the law of total variance:
\begin{equation*}
\text{Var}(X) = E[\text{Var}(X|M)] + \text{Var}(E[X|M])
\end{equation*}

First term (expected conditional variance):
\begin{equation*}
E[\text{Var}(X|M)] = 0.7 \times 1 + 0.3 \times 4 = 0.7 + 1.2 = 1.9
\end{equation*}

Second term (variance of conditional means):
\begin{align*}
\text{Var}(E[X|M]) &= E[(E[X|M])^2] - (E[E[X|M]])^2 \\
&= (0.7 \times 10^2 + 0.3 \times 15^2) - 11.5^2 \\
&= (70 + 67.5) - 132.25 = 137.5 - 132.25 = 5.25
\end{align*}

Total variance:
\begin{equation*}
\text{Var}(X) = 1.9 + 5.25 = 7.15
\end{equation*}

\textbf{Key insight:} The variance of a mixture (7.15) is larger than both component variances (1 and 4) because the variance of conditional means contributes significantly when the modes have different means.
\end{examplebox}

\section{Covariance and Correlation}

Understanding the relationship between multiple random variables is fundamental for multivariate analysis, state estimation, and sensor fusion in control systems.

\begin{definitionbox}[title=Definition: Covariance]
The \textbf{covariance} between two random variables $X$ and $Y$ measures their linear association:
\begin{equation}
\text{Cov}(X, Y) = E[(X - E[X])(Y - E[Y])] = E[XY] - E[X]E[Y]
\end{equation}
\end{definitionbox}

\subsubsection{Properties of Covariance}
\begin{enumerate}
    \item $\text{Cov}(X, X) = \text{Var}(X)$
    \item $\text{Cov}(X, Y) = \text{Cov}(Y, X)$ (symmetry)
    \item $\text{Cov}(aX + b, cY + d) = ac \cdot \text{Cov}(X, Y)$ (linearity)
    \item $\text{Cov}(X + Y, Z) = \text{Cov}(X, Z) + \text{Cov}(Y, Z)$ (distributive)
    \item If $X$ and $Y$ are independent, then $\text{Cov}(X, Y) = 0$
\end{enumerate}

\begin{warningbox}[title=Covariance Zero Does Not Imply Independence]
The converse of property 5 is \textbf{not} true. Two random variables can have zero covariance but still be dependent. Covariance only measures \textit{linear} relationships.

\textbf{Example:} Let $X \sim \mathcal{N}(0,1)$ and $Y = X^2$. Then $\text{Cov}(X, Y) = E[X^3] - E[X]E[X^2] = 0 - 0 = 0$, but $X$ and $Y$ are clearly dependent.
\end{warningbox}

\begin{definitionbox}[title=Definition: Correlation Coefficient]
The \textbf{Pearson correlation coefficient} normalizes covariance to a dimensionless measure:
\begin{equation}
\rho_{XY} = \frac{\text{Cov}(X, Y)}{\sigma_X \sigma_Y} = \frac{E[(X-\mu_X)(Y-\mu_Y)]}{\sqrt{\text{Var}(X)\text{Var}(Y)}}
\end{equation}
where $\rho_{XY} \in [-1, 1]$.
\end{definitionbox}

\textbf{Interpretation:}
\begin{itemize}
    \item $\rho = 1$: Perfect positive linear relationship
    \item $\rho = -1$: Perfect negative linear relationship
    \item $\rho = 0$: No linear relationship (uncorrelated)
    \item $|\rho| \approx 1$: Strong linear relationship
\end{itemize}

\subsubsection{Covariance Matrices}

For a random vector $\vecx = [X_1, X_2, \ldots, X_n]^\top$, the covariance matrix captures all pairwise covariances:

\begin{definitionbox}[title=Definition: Covariance Matrix]
The \textbf{covariance matrix} $\boldsymbol{\Sigma}$ of a random vector $\vecx$ is:
\begin{equation}
\boldsymbol{\Sigma} = E[(\vecx - \boldsymbol{\mu})(\vecx - \boldsymbol{\mu})^\top]
\end{equation}
where $\boldsymbol{\mu} = E[\vecx]$. The $(i,j)$ element is $\Sigma_{ij} = \text{Cov}(X_i, X_j)$.
\end{definitionbox}

Explicitly, the covariance matrix is:
\begin{equation}
\boldsymbol{\Sigma} = \begin{bmatrix}
\text{Var}(X_1) & \text{Cov}(X_1, X_2) & \cdots & \text{Cov}(X_1, X_n) \\
\text{Cov}(X_2, X_1) & \text{Var}(X_2) & \cdots & \text{Cov}(X_2, X_n) \\
\vdots & \vdots & \ddots & \vdots \\
\text{Cov}(X_n, X_1) & \text{Cov}(X_n, X_2) & \cdots & \text{Var}(X_n)
\end{bmatrix}
\end{equation}

\begin{keyconceptbox}[title=Properties of Covariance Matrices]
For any covariance matrix $\boldsymbol{\Sigma}$:
\begin{enumerate}
    \item \textbf{Symmetric:} $\boldsymbol{\Sigma} = \boldsymbol{\Sigma}^\top$
    \item \textbf{Positive semi-definite:} $\vecv^\top \boldsymbol{\Sigma} \vecv \geq 0$ for all $\vecv$
    \item \textbf{Diagonal elements are variances:} $\Sigma_{ii} = \text{Var}(X_i) \geq 0$
    \item \textbf{Off-diagonal bounds:} $|\Sigma_{ij}| \leq \sqrt{\Sigma_{ii}\Sigma_{jj}}$ (follows from $|\rho| \leq 1$)
\end{enumerate}
\end{keyconceptbox}

\textbf{Linear transformation of covariance:} If $\vecy = \mA\vecx + \vecb$, then:
\begin{equation}
\boldsymbol{\Sigma}_{\vecy} = \mA \boldsymbol{\Sigma}_{\vecx} \mA^\top
\end{equation}
This result is fundamental in Kalman filtering when propagating uncertainty through system dynamics.

\section{Stochastic Processes}

A \textbf{stochastic process} is a collection of random variables indexed by time: $\{X(t) : t \in T\}$

\subsubsection{White Noise}
Uncorrelated random sequence with:
\begin{itemize}
    \item $E[w(t)] = 0$
    \item $E[w(t)w(\tau)] = \sigma^2\delta(t-\tau)$
\end{itemize}

\subsubsection{Markov Process}
Future states depend only on present state:
\begin{equation}
P(X_{n+1}|X_n, X_{n-1}, \ldots, X_0) = P(X_{n+1}|X_n)
\end{equation}

\section{Estimation Theory}

\subsubsection{Maximum Likelihood Estimation (MLE)}

\begin{definitionbox}[title=Definition: Maximum Likelihood Estimator]
Given observations $\vecx = (x_1, x_2, \ldots, x_n)$ from a distribution $f(x|\theta)$, the \textbf{maximum likelihood estimator} is:
\begin{equation}
\hat{\theta}_{ML} = \arg\max_\theta \prod_{i=1}^n f(x_i|\theta) = \arg\max_\theta \sum_{i=1}^n \log f(x_i|\theta)
\end{equation}
The second form (log-likelihood) is typically more convenient for computation.
\end{definitionbox}

\begin{examplebox}[title=Example: Maximum Likelihood Estimation for Gaussian Parameters]
Given $n$ independent observations $x_1, x_2, \ldots, x_n$ from a Gaussian distribution $\mathcal{N}(\mu, \sigma^2)$, find the MLEs for $\mu$ and $\sigma^2$.

\textbf{Solution:}

\textbf{Step 1: Write the likelihood function.}
\begin{equation*}
L(\mu, \sigma^2) = \prod_{i=1}^n \frac{1}{\sigma\sqrt{2\pi}} \exp\left(-\frac{(x_i-\mu)^2}{2\sigma^2}\right)
\end{equation*}

\textbf{Step 2: Take the log-likelihood.}
\begin{align*}
\ell(\mu, \sigma^2) &= \log L(\mu, \sigma^2) \\
&= -\frac{n}{2}\log(2\pi) - \frac{n}{2}\log(\sigma^2) - \frac{1}{2\sigma^2}\sum_{i=1}^n (x_i - \mu)^2
\end{align*}

\textbf{Step 3: Find the MLE for $\mu$.}
\begin{equation*}
\frac{\partial \ell}{\partial \mu} = \frac{1}{\sigma^2}\sum_{i=1}^n (x_i - \mu) = 0
\end{equation*}

Solving: $\sum_{i=1}^n x_i = n\mu$, so:
\begin{equation*}
\boxed{\hat{\mu}_{ML} = \bar{x} = \frac{1}{n}\sum_{i=1}^n x_i}
\end{equation*}

\textbf{Step 4: Find the MLE for $\sigma^2$.}
\begin{equation*}
\frac{\partial \ell}{\partial \sigma^2} = -\frac{n}{2\sigma^2} + \frac{1}{2(\sigma^2)^2}\sum_{i=1}^n (x_i - \mu)^2 = 0
\end{equation*}

Solving:
\begin{equation*}
\boxed{\hat{\sigma}^2_{ML} = \frac{1}{n}\sum_{i=1}^n (x_i - \bar{x})^2}
\end{equation*}

\textbf{Note:} The MLE for variance is biased (expected value is $\frac{n-1}{n}\sigma^2$). The unbiased estimator uses $n-1$ in the denominator:
\begin{equation*}
s^2 = \frac{1}{n-1}\sum_{i=1}^n (x_i - \bar{x})^2
\end{equation*}
\end{examplebox}

\subsubsection{Bayesian Estimation}
\begin{equation}
p(\theta|\vecx) \propto p(\vecx|\theta)p(\theta)
\end{equation}
(posterior $\propto$ likelihood $\times$ prior)

\begin{examplebox}[title=Example: Bayesian Parameter Estimation]
Suppose we want to estimate the mean $\mu$ of a Gaussian distribution with known variance $\sigma^2 = 1$. We have a prior belief that $\mu \sim \mathcal{N}(\mu_0, \tau_0^2)$ where $\mu_0 = 0$ and $\tau_0^2 = 4$. After observing a single measurement $x = 2.5$, find the posterior distribution.

\textbf{Solution:}

\textbf{Step 1: Identify the components.}
\begin{itemize}
    \item Prior: $p(\mu) \propto \exp\left(-\frac{(\mu - 0)^2}{2 \times 4}\right) = \exp\left(-\frac{\mu^2}{8}\right)$
    \item Likelihood: $p(x|\mu) \propto \exp\left(-\frac{(x-\mu)^2}{2}\right) = \exp\left(-\frac{(2.5-\mu)^2}{2}\right)$
\end{itemize}

\textbf{Step 2: Apply Bayes' theorem.}
\begin{align*}
p(\mu|x) &\propto p(x|\mu) \cdot p(\mu) \\
&\propto \exp\left(-\frac{(2.5-\mu)^2}{2}\right) \cdot \exp\left(-\frac{\mu^2}{8}\right)
\end{align*}

\textbf{Step 3: Combine the exponents.}
\begin{align*}
-\frac{(2.5-\mu)^2}{2} - \frac{\mu^2}{8} &= -\frac{4(2.5-\mu)^2 + \mu^2}{8} \\
&= -\frac{4(6.25 - 5\mu + \mu^2) + \mu^2}{8} \\
&= -\frac{25 - 20\mu + 4\mu^2 + \mu^2}{8} \\
&= -\frac{5\mu^2 - 20\mu + 25}{8}
\end{align*}

\textbf{Step 4: Complete the square.}
\begin{align*}
5\mu^2 - 20\mu + 25 &= 5(\mu^2 - 4\mu) + 25 \\
&= 5(\mu^2 - 4\mu + 4 - 4) + 25 \\
&= 5(\mu - 2)^2 + 5
\end{align*}

So:
\begin{equation*}
p(\mu|x) \propto \exp\left(-\frac{5(\mu-2)^2}{8}\right) = \exp\left(-\frac{(\mu-2)^2}{2 \times 0.8}\right)
\end{equation*}

\textbf{Result:} The posterior is $\mu|x \sim \mathcal{N}(\mu_n, \tau_n^2)$ where:
\begin{align*}
\tau_n^2 &= \frac{1}{\frac{1}{\tau_0^2} + \frac{1}{\sigma^2}} = \frac{1}{\frac{1}{4} + 1} = \frac{4}{5} = 0.8 \\
\mu_n &= \tau_n^2 \left(\frac{\mu_0}{\tau_0^2} + \frac{x}{\sigma^2}\right) = 0.8\left(\frac{0}{4} + \frac{2.5}{1}\right) = 2.0
\end{align*}

\textbf{Interpretation:}
\begin{itemize}
    \item The posterior mean (2.0) is between the prior mean (0) and the observation (2.5), weighted by their respective precisions.
    \item The posterior variance (0.8) is smaller than both the prior variance (4) and the likelihood variance (1), reflecting increased certainty after incorporating data.
    \item This ``Gaussian-Gaussian'' conjugate pair is the foundation of the Kalman filter.
\end{itemize}
\end{examplebox}

\begin{figure}[htbp]
\centering
\begin{tikzpicture}

% --- Top Panel: Sample Data with Confidence Interval ---
\begin{axis}[
    name=top,
    width=13cm, height=5cm,
    xlabel={Sample Index},
    ylabel={Measured Value},
    xmin=0, xmax=21,
    ymin=3.5, ymax=7.5,
    xtick={1,5,10,15,20},
    grid=major,
    grid style={gray!30},
    axis on top,
    major tick length=0.08cm,
    title={\bfseries Sample Data with 95\% Confidence Interval},
    title style={at={(0.5,1.02)}, anchor=south}
]

% True mean (unknown in practice)
\addplot[UofSCCongaree, thick, dashed, domain=0:21, samples=2] {5.5};
\node[UofSCCongaree, font=\footnotesize, right] at (21, 5.5) {True $\mu$};

% Sample mean
\addplot[UofSCGarnet, very thick, domain=0:21, samples=2] {5.42};
\node[UofSCGarnet, font=\footnotesize, right] at (21, 5.42) {$\bar{x} = 5.42$};

% Confidence interval bounds
\addplot[UofSCAtlantic, thick, dashed, domain=0:21, samples=2] {5.42 + 0.58};
\addplot[UofSCAtlantic, thick, dashed, domain=0:21, samples=2] {5.42 - 0.58};

% Shaded confidence region
\addplot[fill=UofSCAtlantic!20, draw=none, domain=0:21, samples=2] {5.42 + 0.58} \closedcycle;
\addplot[fill=white, draw=none, domain=0:21, samples=2] {5.42 - 0.58} \closedcycle;
\addplot[fill=UofSCAtlantic!20, draw=none] coordinates {(0,4.84) (21,4.84) (21,6.0) (0,6.0)} -- cycle;

% Sample data points
\addplot[only marks, mark=*, mark size=2.5pt, UofSCBlack] coordinates {
    (1, 5.2) (2, 5.8) (3, 4.9) (4, 5.6) (5, 5.1)
    (6, 6.2) (7, 5.0) (8, 5.5) (9, 4.8) (10, 5.9)
    (11, 5.3) (12, 5.7) (13, 4.6) (14, 5.4) (15, 6.1)
    (16, 5.2) (17, 5.8) (18, 5.0) (19, 5.6) (20, 5.5)
};

% CI annotation
\draw[UofSCAtlantic, thick, <->] (1.5, 4.84) -- (1.5, 6.0);
\node[UofSCAtlantic, font=\footnotesize, left, align=right] at (1.5, 5.42) {95\% CI\\$[\bar{x} \pm 1.96\frac{s}{\sqrt{n}}]$};

% Upper/lower bound labels
\node[UofSCAtlantic, font=\footnotesize] at (4, 6.2) {Upper: $\bar{x} + t_{\alpha/2}\frac{s}{\sqrt{n}}$};
\node[UofSCAtlantic, font=\footnotesize] at (4, 4.65) {Lower: $\bar{x} - t_{\alpha/2}\frac{s}{\sqrt{n}}$};

\end{axis}

% --- Bottom Panel: Multiple CI Illustration ---
\begin{axis}[
    name=bottom,
    at={(top.south west)},
    anchor=north west,
    yshift=-1.5cm,
    width=13cm, height=5.5cm,
    xlabel={Experiment Number},
    ylabel={Confidence Interval},
    xmin=0, xmax=21,
    ymin=3.8, ymax=7.2,
    xtick={2,4,6,8,10,12,14,16,18,20},
    grid=major,
    grid style={gray!30},
    axis on top,
    major tick length=0.08cm,
    title={\bfseries Repeated Sampling: 95\% of CIs Contain True Mean},
    title style={at={(0.5,1.02)}, anchor=south}
]

% True mean line
\addplot[UofSCGarnet, very thick, domain=0:21, samples=2] {5.5};
\node[UofSCGarnet, font=\bfseries, right] at (21, 5.5) {$\mu = 5.5$};

% Confidence intervals that capture the mean (Atlantic blue)
\foreach \x/\low/\high in {
    1/4.9/6.0, 2/5.1/6.2, 3/4.8/5.9, 5/5.0/6.1, 6/5.2/6.3,
    7/4.7/5.8, 8/5.0/6.0, 9/4.9/6.0, 10/5.1/6.2, 11/5.0/6.1,
    12/4.8/5.9, 13/5.2/6.3, 15/4.9/6.0, 16/5.0/6.1, 17/4.8/5.9,
    18/5.1/6.2, 19/5.0/6.1, 20/4.9/6.0
}{
    \draw[UofSCAtlantic, thick] (\x, \low) -- (\x, \high);
    \draw[UofSCAtlantic, thick] (\x-0.15, \low) -- (\x+0.15, \low);
    \draw[UofSCAtlantic, thick] (\x-0.15, \high) -- (\x+0.15, \high);
    \node[circle, fill=UofSCAtlantic, minimum size=3pt, inner sep=0pt] at (\x, {(\low+\high)/2}) {};
}

% Confidence intervals that miss the mean (Rose - failures)
\foreach \x/\low/\high in {
    4/5.6/6.7, 14/4.0/5.1
}{
    \draw[UofSCRose, thick] (\x, \low) -- (\x, \high);
    \draw[UofSCRose, thick] (\x-0.15, \low) -- (\x+0.15, \low);
    \draw[UofSCRose, thick] (\x-0.15, \high) -- (\x+0.15, \high);
    \node[circle, fill=UofSCRose, minimum size=3pt, inner sep=0pt] at (\x, {(\low+\high)/2}) {};
}

% Legend annotation
\node[UofSCAtlantic, font=\footnotesize, align=left] at (17, 4.3) {Contains $\mu$};
\draw[UofSCAtlantic, thick] (15.5, 4.3) -- (16.3, 4.3);

\node[UofSCRose, font=\footnotesize, align=left] at (17, 4.0) {Misses $\mu$};
\draw[UofSCRose, thick] (15.5, 4.0) -- (16.3, 4.0);

% Annotation for failures
\draw[UofSCRose, thin, ->] (4, 6.9) -- (4, 6.75);
\node[UofSCRose, font=\footnotesize, above] at (4, 6.9) {Miss!};

\draw[UofSCRose, thin, ->] (14, 4.8) -- (14, 5.0);
\node[UofSCRose, font=\footnotesize, below] at (14, 4.7) {Miss!};

\end{axis}

\end{tikzpicture}
\caption{Confidence intervals for parameter estimation. \textbf{Top}: A sample of $n=20$ measurements (black dots) with sample mean $\bar{x}=5.42$ (Garnet line) and 95\% confidence interval (Atlantic blue shaded region). The true population mean $\mu$ (Congaree dashed) falls within the CI. \textbf{Bottom}: Interpretation of confidence level through repeated sampling. Each vertical bar represents a 95\% CI from an independent experiment. By construction, approximately 95\% of intervals (Atlantic blue, 18/20) contain the true mean, while 5\% (Rose, 2/20) do not. This frequentist interpretation is fundamental to uncertainty quantification in control system identification and sensor calibration.}
\label{fig:confidence_intervals}
\end{figure}


\subsubsection{Minimum Mean Square Error (MMSE)}
\begin{equation}
\hat{\theta}_{MMSE} = E[\theta|\vecx]
\end{equation}

\begin{keyconceptbox}[title=Comparison of Estimators]
\begin{center}
\begin{tabular}{lll}
\toprule
\textbf{Estimator} & \textbf{Criterion} & \textbf{Result} \\
\midrule
MLE & Maximize $p(\vecx|\theta)$ & Mode of likelihood \\
MAP & Maximize $p(\theta|\vecx)$ & Mode of posterior \\
MMSE & Minimize $E[(\theta - \hat{\theta})^2|\vecx]$ & Mean of posterior \\
\bottomrule
\end{tabular}
\end{center}
For symmetric, unimodal posteriors (like Gaussians), MAP and MMSE coincide.
\end{keyconceptbox}

\section{Hypothesis Testing Basics}

Hypothesis testing provides a framework for making decisions under uncertainty, which is essential in fault detection, change detection, and quality control in control systems.

\begin{definitionbox}[title=Definition: Hypothesis Test]
A \textbf{hypothesis test} is a procedure for deciding between two competing claims:
\begin{itemize}
    \item $H_0$: The \textbf{null hypothesis} (default assumption, often ``no effect'' or ``no change'')
    \item $H_1$ (or $H_a$): The \textbf{alternative hypothesis} (the claim we're testing for)
\end{itemize}
\end{definitionbox}

\textbf{Example hypotheses in control systems:}
\begin{itemize}
    \item $H_0$: System is operating normally; $H_1$: Fault has occurred
    \item $H_0$: Controller gain equals design value; $H_1$: Gain has changed
    \item $H_0$: Sensor is calibrated; $H_1$: Sensor has drifted
\end{itemize}

\subsubsection{Types of Errors}

\begin{definitionbox}[title=Definition: Type I and Type II Errors]
\begin{center}
\begin{tabular}{l|cc}
& \textbf{$H_0$ True} & \textbf{$H_1$ True} \\
\hline
\textbf{Reject $H_0$} & Type I Error ($\alpha$) & Correct Decision \\
\textbf{Fail to reject $H_0$} & Correct Decision & Type II Error ($\beta$) \\
\end{tabular}
\end{center}
\begin{itemize}
    \item \textbf{Type I Error (False Positive):} Rejecting $H_0$ when it is true. Probability: $\alpha$ (significance level)
    \item \textbf{Type II Error (False Negative):} Failing to reject $H_0$ when $H_1$ is true. Probability: $\beta$
    \item \textbf{Power:} $1 - \beta$ = probability of correctly rejecting $H_0$ when $H_1$ is true
\end{itemize}
\end{definitionbox}

\begin{warningbox}[title=Trade-off Between Error Types]
There is an inherent trade-off between Type I and Type II errors:
\begin{itemize}
    \item Decreasing $\alpha$ (fewer false alarms) typically increases $\beta$ (more missed detections)
    \item The only way to reduce both simultaneously is to collect more data
\end{itemize}
In control applications, the relative costs of each error type should guide the choice of $\alpha$.
\end{warningbox}

\subsubsection{The p-value}

\begin{definitionbox}[title=Definition: p-value]
The \textbf{p-value} is the probability of observing a test statistic as extreme as, or more extreme than, the observed value, assuming the null hypothesis is true:
\begin{equation}
p\text{-value} = P(\text{test statistic} \geq \text{observed value} \;|\; H_0 \text{ true})
\end{equation}
(For two-sided tests, consider both tails.)
\end{definitionbox}

\textbf{Interpretation:}
\begin{itemize}
    \item Small p-value ($< \alpha$): Strong evidence against $H_0$; reject $H_0$
    \item Large p-value ($\geq \alpha$): Insufficient evidence against $H_0$; fail to reject $H_0$
    \item Common significance levels: $\alpha = 0.05$ (5\%), $\alpha = 0.01$ (1\%)
\end{itemize}

\begin{warningbox}[title=Common Misinterpretations of p-values]
The p-value is \textbf{NOT}:
\begin{itemize}
    \item The probability that $H_0$ is true
    \item The probability that $H_1$ is false
    \item The probability of making an error
\end{itemize}
The p-value is the probability of the \textit{data} (or more extreme), given $H_0$ is true---not the probability of $H_0$ given the data.
\end{warningbox}

\subsubsection{Example: Z-test for a Mean}

\textbf{Setup:} Test whether a population mean $\mu$ differs from a hypothesized value $\mu_0$ when the population variance $\sigma^2$ is known.

\textbf{Hypotheses:} $H_0: \mu = \mu_0$ vs. $H_1: \mu \neq \mu_0$ (two-sided)

\textbf{Test statistic:}
\begin{equation}
Z = \frac{\bar{X} - \mu_0}{\sigma/\sqrt{n}}
\end{equation}
Under $H_0$, $Z \sim \mathcal{N}(0, 1)$.

\textbf{Decision rule:} Reject $H_0$ if $|Z| > z_{\alpha/2}$, where $z_{\alpha/2}$ is the critical value from the standard normal distribution.

\begin{keyconceptbox}[title=Connection to Confidence Intervals]
Hypothesis testing and confidence intervals are dual concepts:
\begin{itemize}
    \item A $100(1-\alpha)\%$ confidence interval contains all parameter values that would \textit{not} be rejected at significance level $\alpha$
    \item Rejecting $H_0: \mu = \mu_0$ at level $\alpha$ is equivalent to $\mu_0$ being outside the $(1-\alpha)$ confidence interval
\end{itemize}
\end{keyconceptbox}

% ============================================================================
